\chapter{算術}
    \section{諸定理}
        \subsection{\texorpdfstring{$a,b,c,d\in\naturalNumbers,\;ab=cd,\;\gcd(a,c)=\gcd(b,d)=1 \Rightarrow a=d,b=c$}{a,b,c,dは自然数, ab=cd, gcd(a,c)=gcd(b,d)=1 ⇔ a=d,b=c}}
            \label{互いに素な2組の数と積の関係}
            \begin{proof}
                \quad\newline
                (最初に思い付いた証明)
                \par
                $p$を、$a$の素因数分解に指数$l\geq1$で現れる任意の素数とする。
                $a$と$c$は互いに素だから$c$の素因数分解に於ける$p$の指数は0である。
                さらに$p\nmid b$である。なぜならば、$b$と$d$が互いに素なので、もし$p\mid b$とすると$d$の素因数分解に於ける$p$の指数は0であり、
                $p\nmid cd$となり$ab=cd$に矛盾するからである。
                よって左辺の素因数分解に於ける$p$の指数は$l$であり、前述の通り$a$と$c$は互いに素だから$d$素因数分解に於ける$p$の指数が$l$でなければならない。
                \par
                $a$と$d$の立場を変えて同様に考えると$d$の素因数分解に指数$m\geq1$で現れる任意の素数$q$は$a$の素因数分解に於いて同じ指数$m$で現れることがわかる。
                よって$a=d$である。これより$b=c$が得られる。
                \newline
                \newline
                (簡潔な証明)
                \par
                $ab=cd$より$a\mid cd$。$a$と$c$は互いに素だから$a\mid d$でなくてはならない。
                同様に、$ab=cd$より$d\mid ab$。$d$と$b$は互いに素だから$d\mid a$でなくてはならない。
                以上より$a=d$。これより$b=c$が得られる。
            \end{proof}
    \section{Extended Euclidean Algorithm}
        Extended Euclidean Algorithm (EEA, 拡張されたユークリッドの互除法)とは、$a,b\in\integers$に対して
        $as + bt = \gcd(a,b)$となるような$s,t\in\integers$を求める手法である。
        この手法では次の規則で非負整数列$\{r_i\}$, および整数列$\{q_i\}, \{s_i\}, \{t_i\}$を定める。
        \[ r_0 = a, s_0 = 1, t_0 = 0,\quad r_1 = b, s_1 = 0, t_1 = 1 \]
        \begin{align*}
            r_{i+1} &= r_{i-1} - q_ir_i,\;0 \leq r_{i+1} < |r_i| \\
            s_{i+1} &= s_{i-1} - q_is_i \\
            t_{i+1} &= t_{i-1} - q_it_i
        \end{align*}
        上の式は、$i\geq 1$については$q_i,r_{i+1}$はそれぞれ$r_{i-1}$を$r_i$で割った商と余りであることを言っている。
        $r_2,r_3,\dotsc$は真に単調減少するから、ある$k\in\naturalNumbers$において$r_{k+1}=0$となる。
        この時点で数列の生成を停止する。
        $a$や$b$が負数の場合は$q_i<0$となることもあるが、それは高々$q_2$までで、$q_3$以降は1以上になる。
        なぜならば$r_2$以降が非負だからである。
        上記の数列について次が成り立つ。
        \begin{enumerate}
            \item $\gcd(a,b) = r_k = as_k + bt_k$
            \item $s_i,t_i\;(i=0,\dotsc,k+1)$および$s_i,s_{i+1}$および$t_i,t_{i+1}\;(i=0,\dotsc,k)$は互いに素である
            \item $|s_{k+1}| = |b|/\gcd(a,b),\;|t_{k+1}| = |a|/\gcd(a,b)$
            \item $a,b>0,\;\gcd(a,b) < \min(a,b)$ならば$|s_i| < b/\gcd(a,b),\;|t_i| < a/\gcd(a,b)\;(i=0,\dotsc,k)$
        \end{enumerate}
        項目4は、計算機でEEAを実行する際に、$a$,$b$を表現できるだけの桁数がある数値型を使えば桁溢れしないことを保証する。
        \begin{proof}
            \quad\\
            \url{https://en.wikipedia.org/wiki/Extended_Euclidean_algorithm}を基に若干改変,加筆した。\\
            (項目1)
            \par
            普通のEuclidの互除法を既に知っているものとする。
            $r_k = \gcd(a,b)$は明らか。
            $r_0 = a = as_0 + bt_0,\;r_1 = b = as_1 + bt_1$であり、$r_2 = r_0 - q_1r_1 = as_0 + bt_0 - q_1(as_1 + bt_1) = a(s_0 - q_1s_1) + b(t_0 - q_1t_1) = as_2 + bt_2$が成り立つ。これを続けると$r_i = as_i + bt_i\;(i=0,\dotsc,k+1)$が得られる。
            \par
            項目1の証明はこれで済んだが、$\{s_i\},\{t_i\}$の漸化式がどうやって考案されたのか想像してみる。
            等式$r_i = as_i + bt_i$が$i=0,1,\dotsc$について成り立って欲しいのだから、$as_{i-1} + bt_{i-1} = r_{i-1} = r_{i+1} + q_ir_i = as_{i+1} + bt_{i+1} + q_i(as_i + bt_i) = a(s_{i+1} + q_is_i) + b(t_{i+1} + q_it_i)$とおいて最左辺と最右辺の$a,b$の係数を比較すれば漸化式を得られる。
            \newline
            (項目2)
            \[
                A_i \coloneq \begin{bmatrix}
                    s_i & s_{i+1} \\
                    t_i & t_{i+1}
                \end{bmatrix},\quad
                Q_i \coloneq \begin{bmatrix}
                    0 & 1 \\
                    1 & -q_i
                \end{bmatrix}
            \]
            とすると
            \[
                A_i = A_0 \prod_{j=1}^i Q_j
            \]
            となる。両辺の行列式を考えると
            \[ s_it_{i+1} - s_{i+1}t_i = |A_i| = |A_0|\prod_{j=1}^i|Q_j| = (-1)^i \]
            これとB\'ezoutの等式より定理の主張が従う。
            \newline
            (項目3)
            \par
            $0 = r_{k+1} = as_{k+1} + bt_{k+1}$より$|a||s_{k+1}| = |b||t_{k+1}|$。
            両辺を$\gcd(a,b)$で割って次式を得る。
            \[ \frac{|a|}{\gcd(a,b)}|s_{k+1}| = \frac{|b|}{\gcd(a,b)}|t_{k+1}| \]
            $|a|/\gcd(a,b)$と$|b|/\gcd(a,b)$は互いに素であり、前述の通り$|s_{k+1}|$と$|t_{k+1}|$も互いに素だから\ref{互いに素な2組の数と積の関係}より
            \[ |s_{k+1}| = \frac{|b|}{\gcd(a,b)},\quad |t_{k+1}| = \frac{|a|}{\gcd(a,b)} \]
            特に$0 = as_{k+1} + bt_{k+1}$を満たすのは次の組み合わせである。
            \[ (s_{k+1},t_{k+1}) = \pm\left(\frac{b}{\gcd(a,b)}, -\frac{a}{\gcd(a,b)}\right) \]
            \newline
            (項目4)
            \par
            まず$s_0,s_2,s_4,\dotsc>0,\;s_3,s_5,s_7,\dotsc<0$を示す。
            $s_0$は定義より正である。さらに$s_2 = s_0 - q_1s_1 = 1,\;s_3 = s_1 - q_2s_2 = -q_2$を直接確かめられる。
            これより始めて帰納的に、$l\geq 2$に対して$s_{2l} = s_{2l-2} - q_{2l-1}s_{2l-1} > 0,\;s_{2l-1} = s_{2l-3} - q_{2l-2}s_{2l-2} < 0$が成り立つ。
            \par
            次に数列$\{|s_i|\}\;(i=1,2,\dotsc)$が単調増加であること、特に$|s_2|\leq|s_3|$を除いて他は全て狭義単調増加であることを示す。
            まず$0 = |s_1| < 1 = |s_2|,\;|s_3| = |-q_2| \geq 1 = |s_2|,\;|s_4| = |s_2 - q_3s_3| = |1 + q_2q_3| > |q_2| = |s_3|$を直接確かめられる。
            これより始めて帰納的に、$l\geq 2$に対して$|s_{2l+1}| - |s_{2l}| = -s_{2l+1} - s_{2l} = (q_{2l}-1)s_{2l} - s_{2l-1} > 0,\;|s_{2l+2}| - |s_{2l+1}| = (1-q_{2l+1})s_{2l+1} + s_{2l} > 0$が成り立つから$|s_{2l+1}| > |s_{2l}|\;(l=2,3,\dotsc)$である。
            同様にして$|t_0| < |t_1| \leq |t_2| < |t_3| < \dotsb$が成り立つ。
            これと項目3より定理の主張が成り立つ。
        \end{proof}